\section{Ba thứ được lắp thêm sau 1997}
\label{sec:ch04-cau-noi}

\index{ô nhớ!ba biến thể sau 1997}%
Kiến trúc trong \code{torch.nn.LSTM} hôm nay không phải kiến trúc mục trước vừa
dẫn. Ba bài dưới đây là ba lần \tn{ô nhớ}{memory cell} được sửa lại sau 1997,
và chúng không cùng một loại. Hai bài đầu đến từ việc dùng bản 1997 rồi va phải
một chỗ nó không xử lý được; chỉ bài thứ nhất trong hai bài ấy nằm trong
\code{torch.nn.LSTM}. Bài thứ ba đến từ hướng khác hẳn.

\begin{bridge}{Gers, Schmidhuber và Cummins 2000: cổng quên}
Bản 1997 giả định chuỗi có điểm kết thúc, và ở đó trạng thái được đặt lại về
không. Bỏ giả định ấy đi, cho mạng đọc một dòng dữ liệu chảy liên tục không
được cắt sẵn thành đoạn, thì phương trình~\eqref{eq:ch04-cap-nhat-c} thành một
phép cộng dồn không bao giờ dừng: \(\vect{c}_t\) lớn dần không giới hạn,
\(g_{\mathrm{out}}\) bão hoà, đạo hàm của nó về không, và ô nhớ chết theo cách
mà chính CEC gây ra.

Cách sửa là bỏ hằng số 1.0 và thay bằng một cổng thứ ba,
\tn{cổng quên}{forget gate} \(\vect{f}_t\), tính giống hai cổng kia:

\[
  \vect{c}_t = \vect{f}_t \odot \vect{c}_{t-1}
             + \vect{i}_t \odot g_{\mathrm{in}}(\vect{a}^{c}_t).
\]

Cái được là ô nhớ tự học lúc nào nên xoá chính nó. Cái mất thì đọc thẳng ra từ
mục~\ref{sec:ch04-dao-ham}: \(\jac{\vect{c}_t}{\vect{c}_{t-1}}\) không còn là
\(\mat{I}\) mà là \(\diag(\vect{f}_t)\), nên vòng quay chỉ không đổi khi cổng
còn mở. Bản 2000 đánh đổi một bảo đảm lấy một khả năng, và đó là bản mà mọi thư
viện hôm nay cài đặt.

Neural Computation 12(10):2451-2471, năm 2000~\autocite{gers2000forget}.
\end{bridge}

\begin{bridge}{Gers, Schraudolph và Schmidhuber 2002: peephole}
Chỗ hở lần này nằm ở chính chỗ mục~\ref{sec:ch04-o-nho} vừa mô tả: cổng nhận
kết nối từ output của các \tn{ô nhớ}{memory cell}, không từ trạng thái. Mà
output là
\(\vect{o}_t \odot g_{\mathrm{out}}(\vect{c}_t)\), nên khi
\tn{cổng ra}{output gate} đóng thì output gần bằng không và không cổng nào còn
nhìn thấy gì. Đóng cổng ra là làm mù cả ba cổng cùng lúc.

Peephole là ba nhóm trọng số nối thẳng từ trạng thái sang các cổng. Thứ tự thời
gian trong đó là chỗ dễ chép sai: \tn{cổng vào}{input gate} và cổng quên nhìn
\(\vect{c}_{t-1}\), còn cổng ra nhìn \(\vect{c}_t\), tức trạng thái đã cập
nhật. Vì thế bài phải chia một bước thành hai pha, cập nhật cổng vào, cổng quên
và trạng thái trước, rồi mới tới cổng ra và output.

Kết quả tiêu đề: mạng phân biệt được chuỗi xung cách nhau 50 bước với chuỗi
cách nhau 49 bước, không cần một ví dụ huấn luyện ngắn nào. Chi phí là ba trọng
số mỗi cell.

Journal of Machine Learning Research 3:115-143, năm
2002~\autocite{gers2002timing}.
\end{bridge}

\begin{bridge}{Cho và cộng sự 2014: hai cổng thay vì ba}
Bài này đến từ hướng khác. Nó không sửa lỗi của LSTM; nó hỏi liệu có cần nhiều
bộ phận đến thế không. Đơn vị họ đề xuất bỏ hẳn trạng thái tách rời và bỏ cổng
ra, giữ hai cổng: reset gate \(\vect{r}_t\) và update gate \(\vect{z}_t\).

\[
  \vect{h}_t = \vect{z}_t \odot \vect{h}_{t-1}
             + (1 - \vect{z}_t) \odot \tilde{\vect{h}}_t,
  \qquad
  \tilde{\vect{h}}_t = \phi\!\left(
    \mat{W}\vect{x}_t + \mat{U}(\vect{r}_t \odot \vect{h}_{t-1})
  \right).
\]

Một cổng duy nhất điều khiển cả giữ lẫn ghi, vì hai hệ số cộng lại bằng 1. Đó
là lý do nó kém linh hoạt hơn: LSTM giữ và ghi độc lập được, đơn vị này thì
không. Về số tham số thì nó ngang bản 1997 và ít hơn bản 2000 đúng một khối,
như bảng ở mục~\ref{sec:ch04-con-lai-gi} đếm ra.

Hai chỗ về tên gọi. Bài \emph{không} dùng chữ \enquote{GRU} và cũng không dùng
\enquote{gated recurrent unit}; nó gọi đây là \enquote{a new type of hidden
unit}, và cái tên viết tắt đến sau. Còn cước chú của chính bài viết rằng LSTM
có \enquote{a memory cell and four gating units}, đếm rộng tay hơn cách đếm
thông thường, vì bản 2000 có ba cổng cộng cell input. Con số bốn hay được lấy
ra từ chỗ này.

EMNLP 2014, trang 1724-1734~\autocite{cho2014encdec}.
\end{bridge}

\subsection*{Vậy \code{torch.nn.LSTM} là bản nào}

\index{ô nhớ!bản trong PyTorch}%
Là bản 2000 có cổng quên, không có peephole, và một chỗ nữa đáng biết: nó giữ
hàm nén trên trạng thái, tức \(\vect{h}_t = \vect{o}_t \odot \tanh(\vect{c}_t)\),
với tanh thật chứ không phải \(g_{\mathrm{out}}\) của bài 1997. Riêng bài 2002
thì bỏ hẳn hàm ấy đi, nói rõ là vì không có bằng chứng thực nghiệm cho thấy nó
cần thiết. Ba bản, ba lựa chọn khác nhau ở đúng một chỗ trong một phương trình.

Nếu bạn muốn một câu tóm tắt để nhớ: cái người ta gọi là LSTM là kiến trúc
1997 với cổng quên của năm 2000 lắp vào, chạy trên gradient đầy đủ mà framework
tính hộ thay vì phép cắt mà bài 1997 đã thiết kế cho nó.
